% HD-R20. GPT-6 Astra Pro /root; inherited effort not exposed.
% 21 November 2026. Actual rank-four generators on two-dimensional blocks.
\section{HD-R20: two-dimensional blocks for the rank-four normal subgroup}
\label{sec:two-block-rank-four}

The complex block-dimension-two branch of the absolutely
irreducible normal rank-four subgroup supplies no missing
sixfold determinant component. Retain the original finite
coefficient group \(G\subset GL_m(K)\), exact imaginary
quadratic character field, signature \((3,3)\), primitive
Jacobian corner and period rank nine. Its actual subgroup
\(N=N_4\), generated by all moving-rank-four elements,
is absolutely irreducible in this branch and preserves
\(r\) complex blocks \(V_i\) of dimension two.
The preceding reduction gives \(m=2r\ge21\), hence
\(r\ge11\). We construct actual descended coefficient
blocks and their original-corner maps.

\paragraph{1. The actual block permutation group.}
For a block cycle of length \(\ell\), successive
transport identifies its fixed vectors with those of its
return map on a two-dimensional block. Its moving rank
is at least \(2(\ell-1)\). Thus the permutation of
an actual rank-four generator satisfies
\[
 \sum_{\ell\ge2}(\ell-1)\le2,
 \qquad\text{its support has size at most four}.          \tag{R20.1}
\]
Its nonidentity cycle type is a transposition, a
three-cycle or two disjoint transpositions. The block
permutation image \(P\) is transitive by irreducibility.

First consider an imprimitive \(P\). Sum the original
two-spaces over a nontrivial permutation block of size
\(a\ge2\). This is an actual complex block system
of dimension \(2a\). Some original generator permutes
these new blocks, since the whole group is transitive
on them. Its moving rank is at least \(2a\), so
\(a=2\). The resulting four-dimensional block system
is covered by the full characteristic-kernel and
original-field projector proof of HD-S50. Its actual
rank-four quotient source gives a supplied determinant
there. This step uses the coarser projectors constructed
in HD-S50 and does not presume descent of the starting
two-dimensional projectors.

Now let \(P\) be primitive. Some generating permutation
is nonidentity, so its minimum support is two, three
or four. The elementary transposition and three-cycle
arguments of HD-S48 give \(S_r\) or a group containing
\(A_r\) in the first two cases. The exact class-four
bound in Manning, \emph{The order of primitive groups. III},
section38, printed p.142, gives degree at most eight
in the third case, excluded by \(r\ge11\).
Thus
\[
                         P=A_r\quad\hbox{or}\quad S_r.
                                                               \tag{R20.2}
\]
The original small-class statement and the elementary
connected-support proof were read and checked in HD-S48;
no new classification premise is introduced here.

Let \(T\) be the actual kernel on these blocks.
Let \(H_i\) be the stabilizer of \(V_i\) in \(N\),
and let \(L_i\subset GL(V_i)\) be the coordinate
image of \(T\). The \(L_i\) are conjugate; write
\(L\) for one of them. The \(H_i\)-representation
on \(V_i\) is irreducible: a proper invariant subspace
would, by summing its transports over the \(N/H_i\)
cosets, give a proper \(N\)-invariant subspace.
The exact quotient is
\[
                H_i/T=P_i=A_{r-1}\quad\hbox{or}\quad S_{r-1}.
                                                               \tag{R20.3}
\]

\paragraph{2. The local kernel is irreducible.}
If \(L\) were scalar, the projective image of \(H_i\)
on its actual two-space would factor through (R20.3).
The finite subgroups of \(PGL_2(\mathbb C)\) are
cyclic, dihedral, \(A_4,S_4,A_5\), in the exact
published form read in HD-S47 (Nguyen--van der Put--Top,
Theorem2, printed p.288). A homomorphism from
\(A_{r-1}\), \(r-1\ge10\), to any of these groups
is trivial. A nontrivial one would be injective by
simplicity; the three polyhedral orders are at most60,
while a dihedral subgroup cannot contain a nonabelian
perfect subgroup, since its derived group is cyclic.
The same is immediate for a cyclic group.
Thus the projective image of \(H_i\) is trivial or
of order two. Its actual linear image is generated
by scalars and at most one finite-order matrix.
These commute and have a common eigenline, contradicting
irreducibility of \(H_i\) on its two-space.

If \(L\) were reducible and nonscalar, finite averaging
would decompose its two-space into two distinct character
lines. The normalizing \(H_i\) must interchange these
lines, since preserving both would contradict its
irreducibility. The interchange homomorphism
\(H_i\to S_2\) kills \(T\), which acts diagonally,
and hence factors through (R20.3). It is trivial for
\(A_{r-1}\), a contradiction. For \(S_{r-1}\)
it is exactly the sign character, since its derived
subgroup is \(A_{r-1}\).

In this last case \(P=S_r\). An original generator
with odd block permutation has, by (R20.1), exactly
one transposition. On each of its \(r-2\) fixed
blocks it belongs to the local \(H_i\), with odd
image in \(S_{r-1}\), and therefore interchanges
the two local character lines. Such a two-by-two
off-diagonal operator fixes at most one dimension:
its fixed-vector equations determine one coordinate
from the other. The two moved blocks contribute at
least two more moving dimensions. Thus this actual
generator has moving rank at least
\[
                              (r-2)+2=r>4,              \tag{R20.4}
\]
a contradiction. Every rank-four generator therefore
has even block permutation, contradicting generation
of \(P=S_r\). This exhausts the reducible cases and
proves that \(L\) is irreducible.

\paragraph{3. The intrinsic kernel without an assumed splitting.}
If \(L\) is solvable, the original embedding
\(T\subset\prod_i L_i\) makes \(T\) solvable.
Conversely every solvable normal subgroup of \(N\)
maps trivially to (R20.2), whose solvable radical is
trivial. Therefore \(T=\operatorname{Rad}(N)\), and
\(T\) is characteristic.

In the nonsolvable case the same finite \(PGL_2\)
classification gives projective coordinate images
\(A_5\). Put
\[
 Z_0=T\cap\prod_i\mathbb C^*I_{V_i},\qquad
                   T/Z_0\subset A_5^r.                 \tag{R20.5}
\]
This is a subdirect product: each coordinate projection
is onto. The subdirect-product argument used in HD-S47
identifies it with a product of diagonal simple factors
on disjoint sets of coordinates. Explicitly, the
two-coordinate image is either the full product or
the graph of an isomorphism, by its two normal kernels
and simplicity. In the graph case the isomorphism is
the literal map \(\pi_j\pi_i^{-1}\) on the common
quotient by their equal kernels. For the induction,
adjoin one coordinate to an already obtained product
of simple factors. The image in the new coordinate
of the old projection kernel is normal, hence is
either all of \(A_5\) or trivial. The first case
adds an independent factor. In the second case the
new map factors through the old direct product;
the commuting factor images are normal in \(A_5\),
so precisely one is onto and all others are trivial.
This gives the graph over that factor. Thus the
product is exactly over the equivalence classes of
equal coordinate kernels. Conjugation
by \(N\) permutes this equivalence relation through
the two-transitive group (R20.2). The only possibilities
are the full product or one diagonal copy of \(A_5\).

The diagonal case is excluded by the original generators.
Every rank-four generator is literally the identity
on at least one fixed block. Indeed its permutation
support has size at most four by (R20.1); among the
other at least seven blocks, at most four can have a
nonidentity diagonal action, since each would contribute
at least one to its total moving rank. Choose such an
identity block for a generator \(x\).
The diagonal group \(T/Z_0\) maps isomorphically to
its projective coordinate image on that block, so
conjugation by \(x\) is the identity on \(T/Z_0\).
Thus \([x,T]\subset Z_0\), and \(Z_0\) is central
in \(T\). Writing \(xtx^{-1}=tz_t\) shows directly
that \(t\mapsto z_t\) is a homomorphism to \(Z_0\).
It kills \(T'\), so \(x\) centralizes \(T'\).
All rank-four generators do; hence \(N\) centralizes
\(T'\). Since \(T'\subset N\), it lies in the scalar
centre of the irreducible \(N\). Its projective image
would be abelian, whereas the image of \(T'\) in
the perfect \(T/Z_0\simeq A_5\) is all of \(A_5\).
This contradiction excludes the diagonal case.

Consequently \(T/Z_0=A_5^r\). The subgroup \(Z_0\)
is abelian and normal in \(N\), and any solvable
normal subgroup lies in \(T\) and has trivial image
in \(A_5^r\). Thus
\(Z_0=\operatorname{Rad}(N)\), an intrinsic subgroup.
In \(\overline N=N/Z_0\), the normal subgroup
\(M=T/Z_0=A_5^r\) is minimal: its simple factors
are permuted transitively. Its centralizer is trivial.
An element centralizing it must fix each simple factor,
so its block permutation is the identity and it lies
in \(M\); the centre of \(M\) is trivial. Any
distinct minimal normal subgroup would commute with
\(M\), which is impossible. Hence \(M\) is the
unique minimal normal subgroup and is characteristic
in \(\overline N\). Its full inverse image \(T\)
is characteristic in \(N\). No complement to \(T\)
has been assumed or constructed in this argument.

\paragraph{4. The original-field projectors.}
The original blocks \(V_i\) are now irreducible
\(T\)-modules. Their isomorphism relation is invariant
under (R20.2), hence they are all equivalent or
pairwise inequivalent. In the first case the literal
evaluation map is
\[
 W\otimes\operatorname{Hom}_T(W,V_{\mathbb C})
       \xrightarrow{\sim}V_{\mathbb C},\qquad
                     w\otimes f\longmapsto f(w),         \tag{R20.6}
\]
with dimensions two and \(r\). Its image algebra
is \(\operatorname{End}(W)\otimes I\); conjugated
matrix units implement the first factor of each
normalizer, and the remaining commutant acts on
the second. Thus (R20.6) is an actual tensor structure
for \(N\), excluded by HD-R17 since \(2r>16\).

The types are therefore distinct. With their characters
\(\chi_i\), the original projectors over a cyclotomic
splitting field \(E=K(\zeta_{\exp T})\) are
\[
 p_i=\frac2{|T|}\sum_{t\in T}\chi_i(t^{-1})\rho(t),
 \qquad p_ip_j=\delta_{ij}p_i,\quad\sum_i p_i=I,
 \quad\operatorname{rank}p_i=2.                          \tag{R20.7}
\]
They are exactly the primitive central projectors of
the actual matrix image \(E[T]\). Galois permutes
them and commutes with the original \(K\)-matrices
of \(N\), hence with (R20.2) on their labels.
That permutation centralizer is trivial: the stabilizer
of a label in \(A_r\) fixes only that label.
Therefore every \(p_i\) belongs to \(M_m(K)\).
Since \(T\) is characteristic in the original
\(N\triangleleft G\), the original \(G\) permutes
these same projectors and the descended spaces
\(U_i=p_iK^m\). Distinct irreducible \(T\)-types
make these spaces orthogonal for the original positive
invariant Hermitian form \(h\).

\paragraph{5. Actual quotient curves and the prescribed primitive corner.}
Set \(J=\operatorname{Stab}_G(U_1)\), let
\(\sigma\) be its actual representation on \(U_1\),
and let \(J_0=\ker\sigma\). The exact induction
and inverse maps are
\[
 F_v(x)=p_1\rho(x)v,\qquad F_v(jx)=\sigma(j)F_v(x),
 \qquad F\longmapsto\sum_{J\backslash G}\rho(x)^{-1}F(x).
                                                               \tag{R20.8}
\]
The inverse identities follow by summing the mutually
orthogonal block projections. Right monodromy and
left deck actions use the actual label inversion
\(Jt^{-1}\leftrightarrow tJ\). On an inertia cycle,
the fixed-section equations are precisely the fixed
equations of its actual return matrix, so (R20.8)
also identifies invariant stalks and first parabolic
cohomology. Its paired map uses the sum of the
orthogonal block forms. The induced trace formula
gives \(K=\mathbb Q(\operatorname{tr}\rho)\subset
\mathbb Q(\operatorname{tr}\sigma)\subset K\).
The faithful quotient \(J/J_0\) therefore has the
same exact character field.

For the original full cover \(\widetilde C\), form
the actual curves and quotient map
\[
       Y=\widetilde C/J,\qquad C'=\widetilde C/J_0,
       \qquad\alpha:\widetilde C\longrightarrow C'.    \tag{R20.9}
\]
For \(0\ne v\in U_1(K)\), put \(q=h(v,v)\).
The original rational primitive idempotent is literally
the induced rank-two one:
\[
 \begin{aligned}
 e_v^G&=\frac{2m}{|G|q}\sum_{g\in G}
             \operatorname{Re}h(v,\rho(g^{-1})v)g\\
      &=\frac4{|J|q}\sum_{j\in J}
             \operatorname{Re}h(v,\sigma(j^{-1})v)j=e_v^J.
 \end{aligned}                                             \tag{R20.10}
\]
Off-\(J\) coefficients vanish and \(|G|=r|J|\).
Averaging over \(J_0\) identifies this with the same
corner of \(\mathbb Q[J/J_0]\). Quotient pullback
is consequently a surjective isogeny from its
integer-cleared primitive image in \(JC'\) onto
the same image in \(J\widetilde C\), with
\[
 \alpha_*\alpha^*=[|J_0|],\qquad
 (\alpha^*)^*\theta_{J\widetilde C}=|J_0|\theta_{JC'},
 \qquad\frac{2|G|}{mq}=\frac{2|J|}{2q}.                \tag{R20.11}
\]
The last equality retains the real Schur normalization
of the paired factor. For the originally prescribed
line \(Kw\), use \(A(x)=v h(w,x)/h(w,w)\) and
the literal Fourier matrix unit
\[
 b=\frac{2m}{|G|}\sum_g\operatorname{Re}\operatorname{tr}
                (A\rho(g^{-1}))g,\qquad
                b^\dagger b=\frac{q}{h(w,w)}e_w.         \tag{R20.12}
\]
After integer clearing by \(M_0\), it is an actual
\(K\)-linear isogeny with polarization multiplier
\(M_0^2q/h(w,w)\). These positive rational multipliers
change a rank-six Hermitian determinant by a sixth
power, the norm of their cube. The original field
action, restricted polarization and period differential
are thus retained through the displayed maps.

The faithful rank-two coefficient system on \(Y\)
has parabolic dimension six and no invariant vectors.
For its genuine coefficient branch values the exact
Euler equation is
\[
               6=2(2g(Y)-2)+\sum_{j=1}^b c_j,
                       \qquad c_j\ge1.                  \tag{R20.13}
\]
The marked-base dimensions are at most seven for
genus zero, six for genus one, and five for genus
two; larger genus is impossible. Indeed the sums
of codimensions in these three cases are ten, six
and two, respectively. For fixed finite monodromy
labels the covers over a marked base are discrete;
the actual quotient variation forgets coefficient-
trivial branch values. The maps (R20.8)--(R20.12)
are flat after fixing those labels, so the original
period rank is at most seven. This contradicts nine.

The primitive permutation case is therefore excluded.
The imprimitive case already reduces through HD-S50
to a supplied rank-four source. Together these actual
constructions close the complete two-dimensional
complex block alternative for the remaining \(N_4\).
